\documentclass{zhvt-classic}

% 使用本地源流明体（GenRyuMin2 TW）按文件路径加载，不依赖系统字体名
% 版框（内框／界栏／书口侧栏／鱼尾）改用与版心印章同一印泥棕色：
%   印章图片 yshs.png 实色 = RGB(135,68,52) = HTML 874434
\usepackage{xcolor}
\usepackage{tikz}
% sealbrown 双注册：xcolor 供 TikZ/pgf，L3 供 cls 版框
\definecolor{sealbrown}{HTML}{874434}
\ExplSyntaxOn
\color_set:nnn { sealbrown } { HTML } { 874434 }
\ExplSyntaxOff

\zhvtset{
  grid color       = sealbrown,
  main font        = GenRyuMin2TW-SB,
  main font path   = C:/Users/Administrator/texmf/fonts/opentype/public/GenRyuMin2TW/,
jiazhu font      = GenRyuMin2TW-M,
  jiazhu font path = C:/Users/Administrator/texmf/fonts/opentype/public/GenRyuMin2TW/,
  % 页码框容纳 4 字：正常页「九十七」+前置空 1 格=4 行恰好；
  % 100+ 页「一百零一」共 5 行超出，末字「一/二/三」被框裁掉（2026-10-04 实测），故放宽到 5
  shukou page number chars = 5,
  % 书口书名框高从默认 6 字减到 4 字，把省出的 2 字高拨给 chapter 框，
  % 使 6 字书名+卷名放进 9 字高的 chapter 框后不再压线（2026-10-05）
  shukou title chars = 4,
}

\title{史记}[孝文本纪]
\author{汉·司马迁}
\maker{闲叟漫记·重製}


% ===== 索隐述赞样式（2026-11-08 四次定稿 · 愚大人亲手改定参数）=====
% ⚠ 竖排前提：自定义字体家族必须带 RawFeature={vertical:+vert}；
%   simhei.ttf 加 +vhal 会整体加载失败静默回退，只写 +vert；simfang 同理。
\newCJKfontfamily{\shuzanfont}[Path=C:/Windows/Fonts/,RawFeature={vertical:+vert}]{simhei.ttf}
\newCJKfontfamily{\shuzanbody}[Path=C:/Windows/Fonts/,RawFeature={vertical:+vert}]{simfang.ttf}
% 题头：152x30pt 圆角 6pt 棕底白字；baseline (0,2) 使整框落入界行格子不压版框线；
% 白字节点 x 不补偿、y +8.2pt（像素级墨迹实测 dx/dy≈0，L/R 留白约 20pt）。
\newcommand{\shuzanhead}[1]{%
  \begin{tikzpicture}[baseline={(0,2)}]
    \fill[sealbrown, rounded corners=6pt]
      (-76pt,-15pt) rectangle (76pt,15pt);
    \node[text=white, inner sep=0pt, anchor=center] at (0pt,8.2pt)
      {\shuzanfont\fontsize{23}{23}\selectfont #1};
  \end{tikzpicture}%
}
% ★ 修正书口卷名位置：chapter 框起始位置改回 cls 默认值，避免全局影响。
%   只在每个卷首页通过 \juan 命令里的空命令微调该页位置（安全，不会跨页污染）。
\ExplSyntaxOn
\AtBeginDocument{
  \dim_set:Nn \l__zhvt_shukou_chapter_title_start_dim
    { \l__zhvt_shukou_title_height_dim + \l__zhvt_font_size_dim * 2 }
  \dim_set:Nn \l__zhvt_shukou_page_number_start_dim
    { \l__zhvt_shukou_chapter_title_start_dim + \l__zhvt_shukou_chapter_title_height_dim }
}
\ExplSyntaxOff

\begin{document}

% ★★ 版心（书口）排布覆写：把书名放到鱼尾之下（2026-10-04）
%
%   cls 默认版心自上而下：书名（title 静态框，在鱼尾之上）→ 鱼尾 → 篇名（chapter 动态框）→ 页码
%   本模板改为：          鱼尾（位置不动）→ 书名 → 篇名 → 页码分割线 → 页码
%
%   手法：把 cls 放在 title 静态框里的书名，改放进 chapter 动态框（鱼尾之下）；
%         title 静态框留空。书名取 \l__zhvt_title_str，即 \title{...} 的第一参数，
%         所以换书时只改 \title，不必在此重复写书名。
%
%   ⚠ 三条必须遵守的写法（2026-10-04 实测）：
%     ① 汉字必须用 \hbox{...} 包住。flowfram 会把框内容装进一个宽 14.4mm（= 一个 line_sep
%        = 书口栏宽）的纵向 minipage 并按该宽度折行，而整页靠 PDF /Rotate 90 显示、
%        框内文字本身不变向 —— 不包 \hbox 时「围炉夜话」4 字会被逐字折成 4 行，
%        在版面上就成了横排 4 字并横向溢出版框。
%     ② 竖排下**一个 \hbox 就是一整列**。书名与篇名要同处一列，必须放进同一个 \hbox；
%        分成两个 \hbox 时，后一个会被当成新的一列、甩到隔壁栏去
%        （实测「序」跑到书口栏外侧 14.4mm 的正文列上）。中间的 \quad 在框内是横向空白，
%        显示出来正好是书名与篇名之间的纵向间隔。
%     ③ ★ 内容盒要写成 **与框等宽的定宽盒** `\hbox to \l__zhvt_shukou_chapter_title_height_dim{…\hfil}`。
%        实测：chapter 框是把内容盒沿纵向**居中**摆放的，内容宽度一变，书名就跟着漂：
%        自然宽度的 \hbox（书名 38.4mm）落在 x=157.2mm，加篇名后（57.6mm）落在 x=152.4mm，
%        目录页（无篇名）与正文页的书名就差半格。做成等宽盒后，无论有没有篇名，
%        内容固定从框顶起排，实测两页均落在 x=142.8mm，不再漂。
\ExplSyntaxOn
% 把书口 chapter 框（书名+篇名）起始位置下移 1 字高，避免 4 字卷名顶到鱼尾
% ⚠ cls 的 \__zhvt_calculate_dimensions: 在 \AtBeginDocument 里执行，会重算该 dim；
%   必须把覆写挂到 begin document 之后，否则会被二次覆盖（2026-10-05 实测）。
% (旧方案已移到导言区，直接修改 chapter 框起始位置)
% 卷标题命令 #1=书口篇名(鱼尾) #2=正文卷名（含卷序，卷名内空格用 ~）
%
% ★★ 目录页两项修订（2026-10-05，愚大人指出）★★
%   ① 页码起点统一对齐：
%      原状 —— 目录项写成 \par\hspace{4em}卷名\hspace{4em}页码，页码起点 = 卷名实宽 + 4em。
%      卷名实宽随字数变：「卷一~梁惠王上」6 字 +~、「卷七~离娄上」5 字 +~，
%      差 1 个字位 → 卷七~卷十 4 条的页码整体比其余 10 条高 27.21pt（一个字），看着上下不齐。
%      改法 —— 卷名补全角空格（宽 = 1 字 = 27.215pt）补足到 **7 字符位**再写进目录，
%      14 条页码起点遂全部落在同一条竖线上（起点对齐；末字随页码字数自然参差）。
%      ⚠ 只对目录里的副本 \l_tmpc_tl 补，正文卷名仍用 \l_tmpa_tl 原样 —— 正文像素不受影响。
%   ② 卷名整体上移半字位：前导 \hspace 由 4em 改 3.5em。
%      实测原状目录块在版框内偏下（天头留白 139.6pt / 地脚留白 110.7pt），
%      上移 0.5em = 13.6pt 后上下基本均分，视觉居中。
\NewDocumentCommand { \juan } { mm } {
	\clearpage
	\tl_set:Nn \l__zhvt_chapter_title_tl { #1 }
	\tl_set:Nn \l_tmpa_tl { #2 }
	\tl_set_eq:NN \l_tmpc_tl \l_tmpa_tl
	\int_step_inline:nn { 7 - \str_count:n { #2 } }
	  { \tl_put_right:Nn \l_tmpc_tl { 　 } }
	\tl_set:Nx \l_tmpb_tl { \zhnum { page } }
	\tl_put_left:Nn \l_tmpb_tl { \hspace { 4em } }
	\tl_put_left:Nn \l_tmpb_tl { \l_tmpc_tl }
	\tl_put_left:Nn \l_tmpb_tl { \par \hspace { 3.5em } }
	\addtocontents { toc } { \l_tmpb_tl }
	\skip_horizontal:n { \l__zhvt_font_size_dim * 4 }
	\l_tmpa_tl
	\par
	\thispagestyle { empty }
}
\cs_set_protected:Npn \__zhvt_set_shukou_content: {
	% ★ 鱼尾上方留白（title 静态框）嵌入闲叟钤印：
	%   整页 /Rotate 90，故此处把图用 \rotatebox{90} 转正为竖章；
	%   尺寸卡在版心条幅内（书内约 7mm 宽 × 33.8mm 高，留余量不压鱼尾）。
	\setstaticcontents * { title } {
		\vbox_to_ht:nn { \l__zhvt_baseline_skip_dim } {
			\vskip 14.2pt
			\hbox_to_wd:nn { \l__zhvt_shukou_title_height_dim } {
				\hfil
				\rotatebox [ origin = c ] { 90 } { \includegraphics [ width = 7mm , height = 33.8mm ] { yshs.png } }
				\hfil
			}
		}
	}
% ★ 版心（书口）统一改法（2026-11-16 愚大人定）：18pt＋框内居中＋书口文字同版框色
%   ① 原 27.2pt 时 7 字长篇名（如「廉颇蔺相如列传」）末字溢出 chapter 框、压住页码；
%      改 \fontsize{18pt}{18pt}，前后 \hfil 令字串在框内居中（不随篇名字数漂移）。
%   ② \vbox{\vskip 8.2pt} 定量下移，抵消缩字号后栏宽方向 8.2pt 上偏（像素实测 y 中心 543.3 = 原版）。
%      ⚠ \raisebox / \vspace* 在 flowfram 框内无效；换字号须同步换补偿：16pt→10.1、19pt→7.4。
%   ③ chapter 与 page~number 均置 \color{sealbrown}（HTML 874434），与版框／界行同色；正文仍黑。
%   ⚠ 页码字位、钤印、鱼尾、界行、正文一律未动；五部合订本同此口径。
	\setdynamiccontents * { chapter }
	{ \color{sealbrown} \vbox { \vskip 8.2pt \hbox to \l__zhvt_shukou_chapter_title_height_dim
		{ \hfil \fontsize{18pt}{18pt}\selectfont \l__zhvt_title_str \quad \l__zhvt_chapter_title_tl \hfil } } }
	\setdynamiccontents * { page~number } { \color{sealbrown} \quad \zhnum { page } }
}
\__zhvt_set_shukou_content:
\ExplSyntaxOff

\maketitle

% 单篇，无多卷目录，不渲染目录页




% 方案①：仅改 A 类错讹，头部与版式不变

\juan{孝文本纪}{卷十~孝文本纪第十}
\hspace*{2em}孝文皇帝。高祖中子也。高祖十一年春。已破陈豨军。定代地。立为代王。都中都。太后薄氏子。即位十七年。高后八年七月。高后崩。九月。诸吕吕产等欲为乱。以危刘氏。大臣共诛之。谋召立代王。事在吕后语中。\\
\hspace*{2em}丞相陈平。太尉周勃等使人迎代王。代王问左右郎中令张武等。张武等议曰。﹁汉大臣皆故高帝时大将。习兵。多谋诈。此其属意非止此也。特畏高帝。吕太后威耳。今已诛诸吕。新喋血京师。此以迎大王为名。实不可信。愿大王称疾毋往。以观其变。﹂中尉宋昌进曰。﹁群臣之议皆非也。夫秦失其政。诸侯豪桀并起。人人自以为得之者以万数。然卒践天子之位者。刘氏也。天下绝望。一矣。高帝封王子弟。地犬牙相制。此所谓盘石之宗也。天下服其彊。二矣。汉兴。除秦苛政。约法令。施德惠。人人自安。难动摇。三矣。夫以吕太后之严。立诸吕为三王。擅权专制。然而太尉以一节入北军。一呼士皆左袒。为刘氏。叛诸吕。卒以灭之。此乃天授。非人力也。今大臣虽欲为变。百姓弗为使。其党宁能专一邪。方今内有硃虚。东牟之亲。外畏吴。楚。淮南。琅邪。齐。代之彊。方今高帝子独淮南王与大王。大王又长。贤圣仁孝。闻於天下。故大臣因天下之心而欲迎立大王。大王勿疑也。﹂代王报太后计之。犹与未定。卜之龟。卦兆得大横。占曰。﹁大横庚庚。余为天王。夏启以光。﹂代王曰。﹁寡人固已为王矣。又何王。﹂卜人曰。﹁所谓天王者乃天子。﹂於是代王乃遣太后弟薄昭往见绛侯。绛侯等具为昭言所以迎立王意。薄昭还报曰。﹁信矣。毋可疑者。﹂代王乃笑谓宋昌曰。﹁果如公言。﹂乃命宋昌参乘。张武等六人乘传诣长安。至高陵休止。而使宋昌先驰之长安观变。\\
\hspace*{2em}昌至渭桥。丞相以下皆迎。宋昌还报。代王驰至渭桥。群臣拜谒称臣。代王下车拜。太尉勃进曰。﹁愿请间言。﹂宋昌曰。﹁所言公。公言之。所言私。王者不受私。﹂太尉乃跪上天子玺符。代王谢曰。﹁至代邸而议之。﹂遂驰入代邸。群臣从至。丞相陈平。太尉周勃。大将军陈武。御史大夫张苍。宗正刘郢。硃虚侯刘章。东牟侯刘兴居。典客刘揭皆再拜言曰。﹁子弘等皆非孝惠帝子。不当奉宗庙。臣谨请阴安侯列侯顷王后与琅邪王。宗室。大臣。列侯。吏二千石议曰。﹃大王高帝长子。宜为高帝嗣。﹄愿大王即天子位。﹂代王曰。﹁奉高帝宗庙。重事也。寡人不佞。不足以称宗庙。愿请楚王计宜者。寡人不敢当。﹂群臣皆伏固请。代王西乡让者三。南乡让者再。丞相平等皆曰。﹁臣伏计之。大王奉高帝宗庙最宜称。虽天下诸侯万民以为宜。臣等为宗庙社稷计。不敢忽。愿大王幸听臣等。臣谨奉天子玺符再拜上。﹂代王曰。﹁宗室将相王列侯以为莫宜寡人。寡人不敢辞。﹂遂即天子位。\\
\hspace*{2em}群臣以礼次侍。乃使太仆婴与东牟侯兴居清宫。奉天子法驾。迎于代邸。皇帝即日夕入未央宫。乃夜拜宋昌为卫将军。镇抚南北军。以张武为郎中令。行殿中。还坐前殿。於是夜下诏书曰。﹁间者诸吕用事擅权。谋为大逆。欲以危刘氏宗庙。赖将相列侯宗室大臣诛之。皆伏其辜。朕初即位。其赦天下。赐民爵一级。女子百户牛酒。酺五日。﹂\\
\hspace*{2em}孝文皇帝元年十月庚戌。徙立故琅邪王泽为燕王。\\
\hspace*{2em}辛亥。皇帝即阼。谒高庙。右丞相平徙为左丞相。太尉勃为右丞相。大将军灌婴为太尉。诸吕所夺齐楚故地。皆复与之。\\
\hspace*{2em}壬子。遣车骑将军薄昭迎皇太后于代。皇帝曰。﹁吕产自置为相国。吕禄为上将军。擅矫遣灌将军婴将兵击齐。欲代刘氏。婴留荥阳弗击。与诸侯合谋以诛吕氏。吕产欲为不善。丞相陈平与太尉周勃谋夺吕产等军。硃虚侯刘章首先捕吕产等。太尉身率襄平侯通持节承诏入北军。典客刘揭身夺赵王吕禄印。益封太尉勃万户。赐金五千斤。丞相陈平。灌将军婴邑各三千户。金二千斤。硃虚侯刘章。襄平侯通。东牟侯刘兴居邑各二千户。金千斤。封典客揭为阳信侯。赐金千斤。﹂\\
\hspace*{2em}十二月。上曰。﹁法者。治之正也。所以禁暴而率善人也。今犯法已论。而使毋罪之父母妻子同产坐之。及为收帑。朕甚不取。其议之。﹂有司皆曰。﹁民不能自治。故为法以禁之。相坐坐收。所以累其心。使重犯法。所从来远矣。如故便。﹂上曰。﹁朕闻法正则民悫。罪当则民从。且夫牧民而导之善者。吏也。其既不能导。又以不正之法罪之。是反害於民为暴者也。何以禁之。朕未见其便。其孰计之。﹂有司皆曰。﹁陛下加大惠。德甚盛。非臣等所及也。请奉诏书。除收帑诸相坐律令。﹂\\
\hspace*{2em}正月。有司言曰。﹁蚤建太子。所以尊宗庙。请立太子。﹂上曰。﹁朕既不德。上帝神明未歆享。天下人民未有嗛志。今纵不能博求天下贤圣有德之人而禅天下焉。而曰豫建太子。是重吾不德也。谓天下何。其安之。﹂有司曰。﹁豫建太子。所以重宗庙社稷。不忘天下也。﹂上曰。﹁楚王。季父也。春秋高。阅天下之义理多矣。明於国家之大体。吴王於朕。兄也。惠仁以好德。淮南王。弟也。秉德以陪朕。岂为不豫哉。诸侯王宗室昆弟有功臣。多贤及有德义者。若举有德以陪朕之不能终。是社稷之灵。天下之福也。今不选举焉。而曰必子。人其以朕为忘贤有德者而专於子。非所以忧天下也。朕甚不取也。﹂有司皆固请曰。﹁古者殷周有国。治安皆千馀岁。古之有天下者莫长焉。用此道也。立嗣必子。所从来远矣。高帝亲率士大夫。始平天下。建诸侯。为帝者太祖。诸侯王及列侯始受国者皆亦为其国祖。子孙继嗣。世世弗绝。天下之大义也。故高帝设之以抚海内。今释宜建而更选於诸侯及宗室。非高帝之志也。更议不宜。子某最长。纯厚慈仁。请建以为太子。﹂上乃许之。因赐天下民当代父後者爵各一级封将军薄昭为轵侯。\\
\hspace*{2em}三月。有司请立皇后。薄太后曰。﹁诸侯皆同姓。立太子母为皇后。﹂皇后姓窦氏。上为立后故。赐天下鳏寡孤独穷困及年八十已上孤兒九岁已下布帛米肉各有数。上从代来。初即位。施德惠天下。填抚诸侯四夷皆洽驩。乃循从代来功臣。上曰。﹁方大臣之诛诸吕迎朕。朕狐疑。皆止朕。唯中尉宋昌劝朕。朕以得保奉宗庙。已尊昌为卫将军。其封昌为壮武侯。诸从朕六人。官皆至九卿。﹂\\
\hspace*{2em}上曰。﹁列侯从高帝入蜀。汉中者六十八人皆益封各三百户。故吏二千石以上从高帝颍川守尊等十人食邑六百户。淮阳守申徒嘉等十人五百户。卫尉定等十人四百户。封淮南王舅父赵兼为周阳侯。齐王舅父驷钧为清郭侯。﹂秋。封故常山丞相蔡兼为樊侯。\\
\hspace*{2em}人或说右丞相曰。﹁君本诛诸吕。迎代王。今又矜其功。受上赏。处尊位。祸且及身。﹂右丞相勃乃谢病免罢。左丞相平专为丞相。\\
\hspace*{2em}二年十月。丞相平卒。复以绛侯勃为丞相。上曰。﹁朕闻古者诸侯建国千馀。各守其地。以时入贡。民不劳苦。上下驩欣。靡有遗德。今列侯多居长安。邑远。吏卒给输费苦。而列侯亦无由教驯其民。其令列侯之国。为吏及诏所止者。遣太子。﹂\\
\hspace*{2em}十一月晦。日有食之。十二月望。日又食。上曰。﹁朕闻之。天生蒸民。为之置君以养治之。人主不德。布政不均。则天示之以菑。以诫不治。乃十一月晦。日有食之。適见于天。菑孰大焉。朕获保宗庙。以微眇之身讬于兆民君王之上。天下治乱。在朕一人。唯二三执政犹吾股肱也。朕下不能理育群生。上以累三光之明。其不德大矣。令至。其悉思朕之过失。及知见思之所不及。匄以告朕。及举贤良方正能直言极谏者。以匡朕之不逮。因各饬其任职。务省繇费以便民。朕既不能远德。故忄间然念外人之有非。是以设备未息。今纵不能罢边屯戍。而又饬兵厚卫。其罢卫将军军。太仆见马遗财足。馀皆以给传置。﹂\\
\hspace*{2em}正月。上曰。﹁农。天下之本。其开籍田。朕亲率耕。以给宗庙粢盛。﹂\\
\hspace*{2em}三月。有司请立皇子为诸侯王。上曰。﹁赵幽王幽死。朕甚怜之。已立其长子遂为赵王。遂弟辟彊及齐悼惠王子硃虚侯章。东牟侯兴居有功。可王。﹂乃立赵幽王少子辟彊为河间王。以齐剧郡立硃虚侯为城阳王。立东牟侯为济北王。皇子武为代王。子参为太原王。子揖为梁王。\\
\hspace*{2em}上曰。﹁古之治天下。朝有进善之旌。诽谤之木。所以通治道而来谏者。今法有诽谤妖言之罪。是使众臣不敢尽情。而上无由闻过失也。将何以来远方之贤良。其除之。民或祝诅上以相约结而後相谩。吏以为大逆。其有他言。而吏又以为诽谤。此细民之愚无知抵死。朕甚不取。自今以来。有犯此者勿听治。﹂\\
\hspace*{2em}九月。初与郡国守相为铜虎符。竹使符。\\
\hspace*{2em}三年十月丁酉晦。日有食之。十一月。上曰。﹁前日遣列侯之国。或辞未行。丞相朕之所重。其为朕率列侯之国。﹂绛侯勃免丞相就国。以太尉颍阴侯婴为丞相。罢太尉官。属丞相。四月。城阳王章薨。淮南王长与从者魏敬杀辟阳侯审食其。\\
\hspace*{2em}五月。匈奴入北地。居河南为寇。帝初幸甘泉。六月。帝曰。﹁汉与匈奴约为昆弟。毋使害边境。所以输遗匈奴甚厚。今右贤王离其国。将众居河南降地。非常故。往来近塞。捕杀吏卒。驱保塞蛮夷。令不得居其故。陵轹边吏。入盗。甚敖无道。非约也。其发边吏骑八万五千诣高奴。遣丞相颍阴侯灌婴击匈奴。﹂匈奴去。发中尉材官属卫将军军长安。\\
\hspace*{2em}辛卯。帝自甘泉之高奴。因幸太原。见故群臣。皆赐之。举功行赏。诸民里赐牛酒。复晋阳中都民三岁。留游太原十馀日。\\
\hspace*{2em}济北王兴居闻帝之代。欲往击胡。乃反。发兵欲袭荥阳。於是诏罢丞相兵。遣棘蒲侯陈武为大将军。将十万往击之。祁侯贺为将军。军荥阳。七月辛亥。帝自太原至长安。乃诏有司曰。﹁济北王背德反上。诖误吏民。为大逆。济北吏民兵未至先自定。及以军地邑降者。皆赦之。复官爵。与王兴居去来。亦赦之。﹂八月。破济北军。虏其王。赦济北诸吏民与王反者。\\
\hspace*{2em}六年。有司言淮南王长废先帝法。不听天子诏。居处毋度。出入拟於天子。擅为法令。与棘蒲侯太子奇谋反。遣人使闽越及匈奴。发其兵。欲以危宗庙社稷。群臣议。皆曰﹁长当弃市﹂帝不忍致法於王。赦其罪。废勿王。群臣请处王蜀严道。邛都。帝许之。长未到处所。行病死。上怜之。後十六年。追尊淮南王长谥为厉王。立其子三人为淮南王。衡山王。庐江王。\\
\hspace*{2em}十三年夏。上曰。﹁盖闻天道祸自怨起而福繇德兴。百官之非。宜由朕躬。今祕祝之官移过于下。以彰吾之不德。朕甚不取。其除之。﹂\\
\hspace*{2em}五月。齐太仓令淳于公有罪当刑。诏狱逮徙系长安。太仓公无男。有女五人。太仓公将行会逮。骂其女曰。﹁生子不生男。有缓急非有益也。﹂其少女缇萦自伤泣。乃随其父至长安。上书曰。﹁妾父为吏。齐中皆称其廉平。今坐法当刑。妾伤夫死者不可复生。刑者不可复属。虽复欲改过自新。其道无由也。妾愿没入为官婢。赎父刑罪。使得自新。﹂书奏天子。天子怜悲其意。乃下诏曰。﹁盖闻有虞氏之时。画衣冠异章服以为僇。而民不犯。何则。至治也。今法有肉刑三。而奸不止。其咎安在。非乃朕德薄而教不明欤。吾甚自愧。故夫驯道不纯而愚民陷焉。诗曰﹃恺悌君子。民之父母﹄。今人有过。教未施而刑加焉。或欲改行为善而道毋由也。朕甚怜之。夫刑至断支体。刻肌肤。终身不息。何其楚痛而不德也。岂称为民父母之意哉。其除肉刑。﹂\\
\hspace*{2em}上曰。﹁农。天下之本。务莫大焉。今勤身从事而有租税之赋。是为本末者毋以异。其於劝农之道未备。其除田之租税。﹂\\
\hspace*{2em}十四年冬。匈奴谋入边为寇。攻朝那塞。杀北地都尉卬。上乃遣三将军军陇西。北地。上郡。中尉周舍为卫将军。郎中令张武为车骑将军。军渭北。车千乘。骑卒十万。帝亲自劳军。勒兵申教令。赐军吏卒。帝欲自将击匈奴。群臣谏。皆不听。皇太后固要帝。帝乃止。於是以东阳侯张相如为大将军。成侯赤为内史。栾布为将军。击匈奴。匈奴遁走。\\
\hspace*{2em}春。上曰。﹁朕获执牺牲珪币以事上帝宗庙。十四年于今。历日长。以不敏不明而久抚临天下。朕甚自愧。其广增诸祀墠场珪币。昔先王远施不求其报。望祀不祈其福。右贤左戚。先民後己。至明之极也。今吾闻祠官祝釐。皆归福朕躬。不为百姓。朕甚愧之。夫以朕不德。而躬享独美其福。百姓不与焉。是重吾不德。其令祠官致敬。毋有所祈。﹂\\
\hspace*{2em}是时北平侯张苍为丞相。方明律历。鲁人公孙臣上书陈终始传五德事。言方今土德时。土德应黄龙见。当改正朔服色制度。天子下其事与丞相议。丞相推以为今水德。始明正十月上黑事。以为其言非是。请罢之。\\
\hspace*{2em}十五年。黄龙见成纪。天子乃复召鲁公孙臣。以为博士。申明土德事。於是上乃下诏曰。﹁有异物之神见于成纪。无害於民。岁以有年。朕亲郊祀上帝诸神。礼官议。毋讳以劳朕。﹂有司礼官皆曰。﹁古者天子夏躬亲礼祀上帝於郊。故曰郊。﹂於是天子始幸雍。郊见五帝。以孟夏四月答礼焉。赵人新垣平以望气见。因说上设立渭阳五庙。欲出周鼎。当有玉英见。\\
\hspace*{2em}十六年。上亲郊见渭阳五帝庙。亦以夏答礼而尚赤。\\
\hspace*{2em}十七年。得玉杯。刻曰﹁人主延寿﹂。於是天子始更为元年。令天下大酺。其岁。新垣平事觉。夷三族。\\
\hspace*{2em}後二年。上曰。﹁朕既不明。不能远德。是以使方外之国或不宁息。夫四荒之外不安其生。封畿之内勤劳不处。二者之咎。皆自於朕之德薄而不能远达也。间者累年。匈奴并暴边境。多杀吏民。边臣兵吏又不能谕吾内志。以重吾不德也。夫久结难连兵。中外之国将何以自宁。今朕夙兴夜寐。勤劳天下。忧苦万民。为之怛惕不安。未尝一日忘於心。故遣使者冠盖相望。结轶於道。以谕朕意於单于。今单于反古之道。计社稷之安。便万民之利。亲与朕俱弃细过。偕之大道。结兄弟之义。以全天下元元之民。和亲已定。始于今年。﹂\\
\hspace*{2em}後六年冬。匈奴三万人入上郡。三万人入云中。以中大夫令勉为车骑将军。军飞狐。故楚相苏意为将军。军句注。将军张武屯北地。河内守周亚夫为将军。居细柳。宗正刘礼为将军。居霸上。祝兹侯军棘门。以备胡。数月。胡人去。亦罢。\\
\hspace*{2em}天下旱。蝗。帝加惠。令诸侯毋入贡。弛山泽。减诸服御狗马。损郎吏员。发仓庾以振贫民。民得卖爵。\\
\hspace*{2em}孝文帝从代来。即位二十三年。宫室苑囿狗马服御无所增益。有不便。辄弛以利民。尝欲作露台。召匠计之。直百金。上曰。﹁百金中民十家之产。吾奉先帝宫室。常恐羞之。何以台为。﹂上常衣綈衣。所幸慎夫人。令衣不得曳地。帏帐不得文绣。以示敦朴。为天下先。治霸陵皆以瓦器。不得以金银铜锡为饰。不治坟。欲为省。毋烦民。南越王尉佗自立为武帝。然上召贵尉佗兄弟。以德报之。佗遂去帝称臣。与匈奴和亲。匈奴背约入盗。然令边备守。不发兵深入。恶烦苦百姓。吴王诈病不朝。就赐几杖。群臣如袁盎等称说虽切。常假借用之。群臣如张武等受赂遗金钱。觉。上乃发御府金钱赐之。以愧其心。弗下吏。专务以德化民。是以海内殷富。兴於礼义。\\
\hspace*{2em}後七年六月己亥。帝崩於未央宫。遗诏曰。﹁朕闻盖天下万物之萌生。靡不有死。死者天地之理。物之自然者。奚可甚哀。当今之时。世咸嘉生而恶死。厚葬以破业。重服以伤生。吾甚不取。且朕既不德。无以佐百姓。今崩。又使重服久临。以离寒暑之数。哀人之父子。伤长幼之志。损其饮食。绝鬼神之祭祀。以重吾不德也。谓天下何。朕获保宗庙。以眇眇之身讬于天下君王之上。二十有馀年矣。赖天地之灵。社稷之福。方内安宁。靡有兵革。朕既不敏。常畏过行。以羞先帝之遗德。维年之久长。惧于不终。今乃幸以天年。得复供养于高庙。朕之不明与嘉之。其奚哀悲之有。其令天下吏民。令到出临三日。皆释服。毋禁取妇嫁女祠祀饮酒食肉者。自当给丧事服临者。皆无践。绖带无过三寸。毋布车及兵器。毋发民男女哭临宫殿。宫殿中当临者。皆以旦夕各十五举声。礼毕罢。非旦夕临时。禁毋得擅哭。已下。服大红十五日。小红十四日。纤七日。释服。佗不在令中者。皆以此令比率从事。布告天下。使明知朕意。霸陵山川因其故。毋有所改。归夫人以下至少使。﹂令中尉亚夫为车骑将军。属国悍为将屯将军。郎中令武为复土将军。发近县见卒万六千人。发内史卒万五千人。藏郭穿复土属将军武。\\
\hspace*{2em}乙巳。群臣皆顿首上尊号曰孝文皇帝。\\
\hspace*{2em}太子即位于高庙。丁未。袭号曰皇帝。\\
\hspace*{2em}孝景皇帝元年十月。制诏御史。﹁盖闻古者祖有功而宗有德。制礼乐各有由。闻歌者。所以发德也。舞者。所以明功也。高庙酎。奏武德。文始。五行之舞。孝惠庙酎。奏文始。五行之舞。孝文皇帝临天下。通关梁。不异远方。除诽谤。去肉刑。赏赐长老。收恤孤独。以育群生。减嗜欲。不受献。不私其利也。罪人不帑。不诛无罪。除刑。出美人。重绝人之世。朕既不敏。不能识。此皆上古之所不及。而孝文皇帝亲行之。德厚侔天地。利泽施四海。靡不获福焉。明象乎日月。而庙乐不称。朕甚惧焉。其为孝文皇帝庙为昭德之舞。以明休德。然后祖宗之功德著於竹帛。施于万世。永永无穷。朕甚嘉之。其与丞相。列侯。中二千石。礼官具为礼仪奏。﹂丞相臣嘉等言。﹁陛下永思孝道。立昭德之舞以明孝文皇帝之盛德。皆臣嘉等愚所不及。臣谨议。世功莫大於高皇帝。德莫盛於孝文皇帝。高皇庙宜为帝者太祖之庙。孝文皇帝庙宜为帝者太宗之庙。天子宜世世献祖宗之庙。郡国诸侯宜各为孝文皇帝立太宗之庙。诸侯王列侯使者侍祠天子。岁献祖宗之庙。请著之竹帛。宣布天下。﹂制曰。﹁可。﹂\\
\hspace*{2em}太史公曰。孔子言﹁必世然後仁。善人之治国百年。亦可以胜残去杀﹂。诚哉是言。汉兴。至孝文四十有馀载。德至盛也。廪廪乡改正服封禅矣。谦让未成於今。呜呼。岂不仁哉。\\
% ===== 索隐述赞（司马贞《史记索隐》）· 统一样式 =====
% 题头「索隐述赞曰」五字并入反白框；正文每四句一列，定宽盒 \hfil 居列中；
% \fontsize{24}{41} 列距对齐全书界行（121+41k）；\clearpage 使述赞独占一页。
\clearpage
\shuzanhead{索隐述赞曰}\linebreak
{\shuzanbody\fontsize{24}{41}\selectfont
\hbox to \hsize{\hfil 孝文在代。兆遇大横。宋昌建册。绛侯奉迎。\hfil}\linebreak
\hbox to \hsize{\hfil 南面而让。天下归诚。务农先籍。布德偃兵。\hfil}\linebreak
\hbox to \hsize{\hfil 除帑削谤。政简刑清。綈衣率俗。露台罢营。\hfil}\linebreak
\hbox to \hsize{\hfil 法宽张武。狱恤缇萦。霸陵如故。千年颂声。\hfil}%
\par}
\end{document}
